% Generated draft: requires mathematical and source review.
\documentclass[11pt]{article}
\usepackage[margin=25mm]{geometry}
\usepackage{fontspec}
\setmainfont{Latin Modern Roman}
\usepackage{amsmath,amssymb,mathtools,unicode-math}
\setmathfont{Latin Modern Math}
\usepackage{enumitem,needspace,xurl,hyperref}
\setlength{\parindent}{0pt}
\setlength{\parskip}{5pt}
\setlength{\emergencystretch}{3em}
\newcommand{\N}{\mathbb{N}}
\newcommand{\Z}{\mathbb{Z}}
\newcommand{\Q}{\mathbb{Q}}
\newcommand{\R}{\mathbb{R}}
\newcommand{\C}{\mathbb{C}}
\newcommand{\F}{\mathbb{F}}
\newcommand{\A}{\mathbb{A}}
\newcommand{\PP}{\mathbb{P}}
\newcommand{\E}{\mathbb{E}}
\newcommand{\eps}{\varepsilon}
\newcommand{\dd}{\,\mathrm{d}}
\begin{document}
\section*{UnsolvedMath ID: 700603}
\textbf{Source problem code:} catalogue entry
\subsection{Definitions and mathematical statement}
Let \(G\leq\operatorname{Sym}(\Omega)\) be a finite transitive permutation group on a nonempty finite set \(\Omega\). The action is primitive if its only \(G\)-invariant partitions are the partition into singletons and the partition with sole part \(\Omega\). Fix \(\alpha\in\Omega\) and put \(H=G_\alpha=\{g\in G:g\alpha=\alpha\}\). The rank of \(G\) is
\[
 r(G)=|H\backslash\Omega|,
\]
where \(H\backslash\Omega\) denotes the set of \(H\)-orbits on \(\Omega\).

For each \(H\)-orbit \(\Delta\subseteq\Omega\), let \(H^\Delta\leq\operatorname{Sym}(\Delta)\) be the permutation group induced by \(H\) on \(\Delta\). This transitive action is a transitive constituent of the point stabilizer. For \(\beta\in\Delta\), define
\[
 r(H^\Delta)=|H_\beta\backslash\Delta|,
 \qquad H_\beta=\{h\in H:h\beta=\beta\}.
\]
This rank is independent of the choice of \(\beta\in\Delta\). The subrank of \(G\) is
\[
 s(G)=\max_{\Delta\in H\backslash\Omega}r(H^\Delta).
\]
Both \(r(G)\) and \(s(G)\) are independent of the chosen point \(\alpha\). Singleton constituents are included and have rank \(1\).

\textbf{Conjecture of P. Cameron, as posed in Kourovka Notebook Problem 9.70.} There exists a function
\[
 f:\Z_{\geq1}\longrightarrow\Z_{\geq1}
\]
such that every finite primitive permutation group \(G\leq\operatorname{Sym}(\Omega)\), with \(m=s(G)\), satisfies
\[
 r(G)\leq f(m)\qquad\text{or}\qquad |G_\alpha|\leq m.
\]
Thus the bound on rank must depend only on the subrank, uniformly over all such groups and degrees. Equivalently, for each positive integer \(m\), the ranks of finite primitive permutation groups with subrank \(m\) and point stabilizer order greater than \(m\) are uniformly bounded.

The supplied source does not specify whether one-point actions are admitted as primitive. The partition definition above admits them; their rank, subrank, and point stabilizer order are all \(1\), so their inclusion does not change the conjecture.
\subsection{Short English statement}
Cameron's conjecture asks whether the subrank uniformly bounds the rank of a finite primitive permutation group whenever its point stabilizer has order greater than the subrank. Subrank measures the largest rank among the actions of a point stabilizer on its individual orbits. The supplied registry labels this problem open; that status has not been independently verified.
\subsection{Sources}
\begin{itemize}
\item Local UnsolvedMath registry, supplied complete source record: permanent ID \(700603\), exact source code \texttt{catalogue entry}, title \emph{Kourovka Notebook Problem 9.70}. The supplied status is \texttt{open}.
\item \emph{The Kourovka Notebook (Unsolved Problems in Group Theory)}, Problem 9.70, proposed by A. N. Fomin in the 9th issue (1984). The supplied background identifies the 21st-issue edition (2026), arXiv:1401.0300v48. The collection description names E. I. Khukhro and V. D. Mazurov as editors.
\item P. Cameron, \emph{Bulletin of the London Mathematical Society}, volume \(13\), no. \(1\) (1981), pp. 1--22. This is the reference supplied for the conjecture; no article title or external link is supplied.
\end{itemize}
\end{document}
